\chapter{Functionalities and user interface}
\label{chap:functionalities}

This chapter describes FlowSync from the point of view of its users: the principles that guided
the design of the interface, how the pages are organised by role, and the main flows, from an
order to an approved plan and from a reported problem to a new plan. The implementation of the
interface is described in Section~\ref{sec:impl-frontend}.

\section{Design principles}
\label{sec:ui-principles}

\begin{itemize}
  \item \textbf{Decision support, not autopilot.} Every page that shows a plan also shows where
        it comes from (the model or the fallback rule, with its confidence) and what it costs.
        The planner always has the last word, and the buttons that change the state of a plan
        ask for a confirmation or a reason.
  \item \textbf{Alternatives instead of parameters.} The planner rarely knows in advance which
        weight to give to punctuality or overtime. FlowSync therefore shows three concrete
        plans (economic, balanced, fast) together with the classical rules ATC and EDD, and lets
        the planner choose by comparing their indicators.
  \item \textbf{Say why.} When something is not allowed, the interface says why, in words and
        next to the element that caused it: a move on the Gantt chart that violates a
        precedence is undone with the name of the activity that must finish first, a cycle in
        the process graph is rejected with the activities involved, a form error appears under
        its field and receives the focus.
  \item \textbf{One visual system.} All the pages share the same colours, typography and
        components, defined once as design tokens; the colours of the states (to review,
        approved, rejected, superseded) and of the orders on the Gantt chart are the same
        everywhere.
  \item \textbf{Accessible.} Text and controls respect the contrast ratios of WCAG~2.1 AA, every
        control is reachable from the keyboard with a visible focus, interactive targets are at
        least 44 pixels high, and each drag-and-drop action has a keyboard alternative. The
        interface is designed for desktop screens, which is where planners work.
\end{itemize}

\section{Pages and roles}
\label{sec:ui-pages}

After the login, each role sees only its own pages. The planner and the administrator use a
layout with a side menu; the operator uses a simpler layout with no menu, built around the
activities of the day. Table~\ref{tab:pages} lists the pages.

\begin{table}[htpb]
  \centering
  \small
  \begin{tabular}{@{}p{3.4cm}p{2.4cm}p{8cm}@{}}
    \toprule
    \textbf{Page} & \textbf{Roles} & \textbf{Purpose} \\
    \midrule
    Login, registration & all & Access with username and password; request of a new
      account \\
    Dashboard & planner, admin & Indicators of the last weeks, open alerts, utilisation per
      resource, recent requests \\
    Orders & planner, admin & Orders, process graph of the activities, alternative resources
      and durations \\
    New planning & planner & Request in natural language, interpretation, missing data,
      weights \\
    Plans & planner, admin & All the requests with their plans, with search and filters \\
    Plan comparison & planner, admin & The three plans of a request and the comparison with ATC
      and EDD \\
    Gantt and review & planner, admin & A plan on the resource Gantt chart; moves, approval,
      rejection, new version \\
    Reschedule & planner & A new plan after a problem in the running plan \\
    Today's activities & operator & Start, pause, finish; problem reports \\
    Configuration & planner (read), admin & Resources, calendars, skills, domain profiles \\
    AI model & planner, admin & Active models, metrics, fallback chain, version history \\
    Users & admin & Pending registrations to approve with a role, accounts to enable or
      disable \\
    Audit & admin & Who did what, when and with which model; CSV export \\
    \bottomrule
  \end{tabular}
  \caption{Pages of FlowSync}
  \label{tab:pages}
\end{table}

\section{From an order to an approved plan}
\label{sec:ui-main-flow}

Figure~\ref{fig:main-flow} shows the main flow of the planner.

\begin{figure}[htpb]
  \centering
  \begin{tikzpicture}[
      font=\small,
      box/.style={draw, rounded corners, align=center, minimum width=2.3cm,
                   minimum height=1.1cm, fill=gray!8},
      ai/.style={draw, dashed, rounded corners, align=center, minimum width=2.3cm,
                 minimum height=1.1cm},
      arrow/.style={-{Stealth[length=2mm]}, thick},
      node distance=0.45cm]
    \node[box] (orders) {Orders and\\process graph};
    \node[box, right=of orders] (request) {Request\\(text or form)};
    \node[box, right=of request] (check) {Check the\\interpretation};
    \node[ai, right=of check] (compute) {3 plans\\+ ATC, EDD};
    \node[box, below=0.9cm of compute] (compare) {Compare\\the plans};
    \node[box, left=of compare] (gantt) {Review on\\the Gantt};
    \node[box, left=of gantt] (decide) {Approve,\\modify, reject};
    \node[box, left=of decide] (execute) {Execution by\\the operators};
    \draw[arrow] (orders) -- (request);
    \draw[arrow] (request) -- (check);
    \draw[arrow] (check) -- (compute);
    \draw[arrow] (compute) -- (compare);
    \draw[arrow] (compare) -- (gantt);
    \draw[arrow] (gantt) -- (decide);
    \draw[arrow] (decide) -- (execute);
  \end{tikzpicture}
  \caption{Main flow of the planner. The dashed box is computed by the AI service, through the
           backend; all the other steps are decisions of the user.}
  \label{fig:main-flow}
\end{figure}

\paragraph{Orders and process graph.} An order has a customer, a release date, a due date (which
may still be unknown), a priority and the weight of its delay. Its activities are drawn as nodes
of a graph, and a precedence is created by dragging from one activity to another; a precedence
that would close a cycle is refused at once. Each activity lists the skills it requires and one
or more alternatives, each with the resource that can execute it and its duration; when the
duration is not known, it is left empty and estimated by the AI. The same information can be
edited in a table, which is the keyboard alternative to the graph. Changes to the activities and
precedences are kept as a draft and saved together; leaving the page with unsaved changes asks
for a confirmation.

\paragraph{The request.} The planner can describe what is needed in words, for example
\emph{``plan C-101 to C-104 for next week, avoiding overtime''}. A language model turns the text
into the parameters of the request: the orders, the horizon, the hard constraints (due dates,
maintenance windows) and the preferences, which become the weights of the objectives. The page
shows this interpretation, explains why a weight changed, and asks the questions about the
missing data, such as a due date that was never entered. Nothing is computed until the planner
has checked the interpretation and answered the questions.

\paragraph{Waiting for the plans.} The request is computed in the background. The page shows the
stages of the computation (validation of the data, estimation of the missing durations,
construction of the plans, comparison with the rules), which take a few seconds, and then
replaces itself with the result.

\paragraph{Comparing the plans.} The three plans are shown side by side, each with a sentence
that summarises its trade-off (\emph{``no delay, two hours of overtime on Tuesday''}) and its
indicators: total duration, total tardiness, overtime, number of setups, average utilisation and
estimated cost. The plan that best matches the weights of the planner is marked as recommended.
Below, a table compares the plans of the AI with those of the ATC and EDD rules on the same data.

\paragraph{Review on the Gantt chart.} Each plan opens on a resource Gantt chart: one row per
resource, one bar per activity, a colour per order, grey bands where a resource is unavailable,
striped bars where the duration was estimated by the AI. The planner can drag a bar in time or
onto another resource; the move is checked immediately against the precedences, the other
activities of the resource, the unavailability and the resources allowed for the activity, and
the indicators are recomputed. A side panel shows the details of the selected activity, including
the estimate of the AI with its interval, and offers buttons to move it by fifteen minutes or an
hour, or to another resource.

\paragraph{Decision.} The planner approves the plan (the other alternatives of the request are
discarded), rejects it with a reason, or saves the modified plan as a new version with a reason.
The reason ends up in the audit log. A decided plan becomes read-only.

\section{When something goes wrong}
\label{sec:ui-reschedule}

Plans meet reality: a machine breaks down, an order becomes urgent, an activity takes longer
than expected. In FlowSync the problem is reported by the operator who sees it, from the page of
the activities of the day: the operator chooses the kind of problem (breakdown, missing material,
delay, other), the machine involved and, if useful, a note. The report reaches the planner
immediately as an alert on the dashboard.

Since the problem is already known, the planner is not asked to describe it again. The alert
opens the rescheduling page, which shows:
\begin{itemize}
  \item the problem, as it was reported (for example \emph{``Welding machine S2 down, reported
        by M.~Rossi at 10:42: it does not turn on''});
  \item its impact on the running plan: the activities blocked by the problem, those that come
        after them and will slip, the orders involved and their due dates;
  \item a few optional parameters: until when the machine will be down (if the technician gave
        a time), the weights of the objectives, and how much the plan may change, from
        \emph{change as little as possible} to \emph{re-optimise everything}.
\end{itemize}
A single button asks the AI to reschedule. The result is a new request, with three plans and the
comparison with the rules, like any other; activities already started stay where they are. The
current plan remains in execution until the planner approves one of the new plans, which then
replaces it. If the AI cannot find a feasible plan, the page says so and why. When the
rescheduling does not start from a report, for instance from the button on the Gantt chart, the
page first asks what happened.

\section{The operator}
\label{sec:ui-operator}

The operator's page lists the activities of the day assigned to the operator's resources: the
activity in progress, with a timer and large buttons to pause or finish it, the next activities,
and those already finished. An activity cannot be started if a preceding one has not finished
(the page says which one and when it is expected to end), and only one activity can be in
progress at a time. The recorded start and end times are not only a record of the work: they
are the data from which the duration model M1 learns.

\section{Transparency and control}
\label{sec:ui-transparency}

Three pages make the behaviour of the system visible.
\begin{itemize}
  \item The \textbf{dashboard} shows the indicators of the last weeks (punctuality, late orders,
        utilisation, adherence to the plan), the open alerts and the recent requests.
  \item The \textbf{AI model} page shows which versions of M1 and M2 are active and what their
        metrics mean, in plain words (for example, \emph{gap from the optimum: how much worse
        the plans of the AI are than those computed by CP-SAT}); the fallback chain, with how many
        plans of the last thirty days came from the AI and how many from a fallback; and the
        history of the versions. The page is read-only: a model is published only by the
        training pipeline.
  \item The \textbf{audit} page lists who did what and when: plans generated, decisions,
        reschedulings, fallbacks. Each record shows the model version and the hashes of input and
        output, with which the plan can be reproduced, and can be exported as CSV.
\end{itemize}

\section{Users}
\label{sec:ui-users}

Anyone can ask for an account from the registration page, giving name, username, e-mail and
password; the account is created \emph{pending}, without a role. The administrator finds the
pending requests at the top of the users page, with their number also shown next to the menu
entry, and approves each of them by choosing a role, or rejects it. Below, the list of the
accounts lets the administrator change a role or disable and re-enable an account; the
administrator cannot disable or demote their own account, so the system always keeps an
administrator. A user who tries to log in while pending or disabled is told why.

\section{Configuration}
\label{sec:ui-configuration}

The administrator maintains the master data on four tabs: resources (type, skills, calendar,
hourly cost, planned unavailability), calendars (working days, shifts and closures, with the
resulting weekly hours), skills (with the active resources that have each one, and a warning
when no resource has it) and domain profiles (default weights of the objectives and active
constraints). Planners can read these data but not change them.

% TODO (M. Brizi): screenshots of the main pages (dashboard, plan comparison, Gantt, reschedule,
% operator) in Immagini/, once the data of the real backend are available.
